% sample-flowchart.tex
% Standalone or includable TikZ flowchart snippet.

\begin{figure}[htbp]
  \centering
  \begin{tikzpicture}[
    node distance=1.5cm,
    auto,
    block/.style={
      rectangle,
      draw=black!75,
      fill=blue!10,
      thick,
      rounded corners=3pt,
      text width=3.2cm,
      text centered,
      minimum height=1.1cm,
      font=\small\sffamily
    },
    decision/.style={
      diamond,
      draw=black!75,
      fill=yellow!20,
      thick,
      aspect=2,
      text width=2.5cm,
      text centered,
      font=\footnotesize\sffamily
    },
    line/.style={
      draw=black!85,
      -stealth,
      thick
    }
  ]
    \node [block] (init) {Initialisation des données};
    \node [block, below=of init] (process) {Calcul des paramètres et optimisation};
    \node [decision, below=of process] (decide) {Convergence atteinte ?};
    \node [block, right=2cm of decide] (update) {Ajustement du pas de gradient};
    \node [block, below=of decide] (stop) {Génération du rapport et métriques};

    \path [line] (init) -- (process);
    \path [line] (process) -- (decide);
    \path [line] (decide) -- node [near start, right] {Oui} (stop);
    \path [line] (decide) -- node [near start, above] {Non} (update);
    \path [line] (update) |- (process);
  \end{tikzpicture}
  \caption{Schéma algorithmique du processus d'optimisation itérative.}
  \label{fig:flux-optimisation}
\end{figure}

